pdfoutput=1
\pdfoutput=1
% ============================================================
% S3: Theorem 3 -- The Honest Person Theorem (The Honest Person Theorem)
% Supplementary Information for Nature submission:
% "A Fundamental Impossibility in Data Quality"
% ============================================================

\section{Theorem 3: The Honest Person Theorem (The Honest Person Theorem)}
\label{sec:thm3}

\subsection{Statement and Implications}
\label{sec:thm3_statement}

The central discovery of this paper is that distinguishing label noise from
intrinsic sample difficulty is mathematically impossible from observational data
alone.  Theorem~3 formalizes this impossibility.

\begin{theorem}[The Honest Person Theorem (The Honest Person Theorem)]
  \label{thm:unidentifiability}
  For any $K \ge 2$ classification problem, any $M \ge 1$ experts, and any finite
  state space~$\mathcal{S}$, there exist two data-generating processes
  $\mathcal{P}_{\text{noise}}$ and $\mathcal{P}_{\text{hard}}$ such that:
  \begin{enumerate}[label=(\roman*)]
    \item Under $\mathcal{P}_{\text{noise}}$, label errors are caused by
      \textbf{label noise}---the observed label differs from the true label.
    \item Under $\mathcal{P}_{\text{hard}}$, all observed labels equal the true
      labels, but some samples are intrinsically
      \textbf{difficult}---expert errors arise from label ambiguity, not from
      label errors.
    \item The two processes produce identical observable joint distributions
      \textbf{if} the noise rate $\eta_{\text{err}}$ in $\mathcal{P}_{\text{noise}}$
      equals the ambiguity rate $\eta_{\text{amb}}$ in $\mathcal{P}_{\text{hard}}$.
      (Note: only the sufficiency direction is proved---the constructed worlds
      are observationally equivalent when $\eta_{\text{err}}=\eta_{\text{amb}}$;
      necessity, i.e.\ whether observational equivalence forces equality of the
      two rates for \emph{all possible} world pairs, is \textbf{not established}.)
      Under this symmetry condition $\eta_{\text{err}} = \eta_{\text{amb}} = \eta$:
      \[
      \mathcal{P}_{\text{noise}}\bigl(x,\,y,\,\{f_m(x)\}\bigr)
      = \mathcal{P}_{\text{hard}}\bigl(x,\,y,\,\{f_m(x)\}\bigr),
      \qquad \forall\,(x,y,\{f_m\}) \in \mathcal{X}\times\mathcal{Y}\times\mathcal{Y}^M.
      \]
    \item Consequently, for any algorithm $\mathcal{A}$ that maps $n$ i.i.d.\
      observations to per-sample noise flags, the expected error rate on the
      ambiguous subset satisfies
      \[
      \max\bigl(\mathrm{Error}_{\mathcal{P}_{\text{noise}}}(\mathcal{A}),\,
             \mathrm{Error}_{\mathcal{P}_{\text{hard}}}(\mathcal{A})\bigr)
      \ge \frac{\eta\rho}{2},
      \]
      where $\eta$ is the noise rate and $\rho$ is the proportion of the
      ambiguous state.  The bound is strictly positive whenever $\eta>0$ and
      $\rho>0$, proving that the two worlds cannot be perfectly distinguished.
  \end{enumerate}
\end{theorem}

\begin{remark}
  Theorem~3 asserts that the two statements ``this sample's label is wrong''
  and ``this sample is so difficult that experts systematically err on it'' are
  observationally equivalent.  The SCX framework must introduce assumptions
  (A1)--(A6) to break this equivalence; Theorem~3 shows that these assumptions
  are not arbitrary technical conditions but rather the minimal structure
  required for identifiability.  See \S\ref{sec:thm3_minimal} for the minimal
  sufficient set.
\end{remark}

\subsection{Proof for $K=2$ (Binary Classification)}
\label{sec:thm3_proof_K2}

We prove Theorem~3 by explicit construction.  Fix $K=2$, label space
$\mathcal{Y}=\{0,1\}$, $M\ge 1$ experts, and two states $s_1,s_2\in\mathcal{S}$.
Let $\rho\in(0,1)$ and $\eta\in(0,\tfrac12)$ be parameters to be specified.

\subsubsection{World A: Noise-Driven Construction}
\label{sec:thm3_worldA}

In the noise world $\mathcal{P}_{\text{noise}}$:

\begin{itemize}
  \item \textbf{State $s_1$} (``clean but noisy''), with
    $\mathbb{P}(X\in s_1)=\rho$.  All samples have true label $y^*\equiv0$.
    Labels are flipped with probability $\eta_{\text{err}}$:
    \[
    \mathbb{P}(y\neq y^* \mid x\in s_1)=\eta_{\text{err}}.
    \]
    Expert accuracy on clean samples is $1-\varepsilon_1$:
    \[
    \mathbb{P}(f_m(x)=y^* \mid x\in s_1,\text{ clean})=1-\varepsilon_1,
    \qquad \varepsilon_1\in(0,\tfrac12).
    \]
  \item \textbf{State $s_2$} (``clean and easy''), with
    $\mathbb{P}(X\in s_2)=1-\rho$.  True label $y^*\equiv0$, no noise
    ($\eta=0$).  Expert accuracy is $1-\varepsilon_2$ with
    $\varepsilon_1<\varepsilon_2\le\tfrac12$ (the state is genuinely harder
    for experts).
\end{itemize}

Expert predictions depend only on the true label $y^*$ (through training on
clean data) and are conditionally independent given $x$.

\subsubsection{World B: Difficulty-Driven Construction}
\label{sec:thm3_worldB}

In the hardness world $\mathcal{P}_{\text{hard}}$:

\begin{itemize}
  \item \textbf{State $s_1$}, with the same marginal probability $\rho$.
    All observed labels are correct: $y=y^*$ for every sample.  However, the
    true label itself is random:
    \[
    \mathbb{P}(y^*=0 \mid x\in s_1)=1-\eta_{\text{amb}},\qquad
    \mathbb{P}(y^*=1 \mid x\in s_1)=\eta_{\text{amb}}.
    \]
    Experts are biased toward class~0: for every $m$,
    \[
    \mathbb{P}(f_m(x)=0 \mid x\in s_1,\,y^*=0)=1-\varepsilon_1,\qquad
    \mathbb{P}(f_m(x)=0 \mid x\in s_1,\,y^*=1)=1-\varepsilon_1.
    \]
    Thus when $y^*=0$ the expert accuracy is $1-\varepsilon_1$ (good); when
    $y^*=1$ the expert accuracy is $\varepsilon_1$ (systematically wrong).
  \item \textbf{State $s_2$}, identical to World A: $y=y^*\equiv0$, expert
    accuracy $1-\varepsilon_2$.
\end{itemize}

\subsubsection{Observational Equivalence Verification}
\label{sec:thm3_verification}

The observable joint distribution factorises as
\[
\mathcal{P}(x,y,f_1,\dots,f_M)
= \mathcal{P}(x)\cdot\mathcal{P}(y\mid x)\cdot\prod_{m=1}^M\mathcal{P}(f_m\mid x),
\]
using conditional independence of experts given $x$ (Assumptions A1--A2).

\paragraph{Marginal $\mathcal{P}(x)$.}
In both worlds $\mathbb{P}(X\in s_1)=\rho$,
$\mathbb{P}(X\in s_2)=1-\rho$.  Identical.

\paragraph{Conditional label distribution $\mathcal{P}(y\mid x)$.}
In World A (state $s_1$):
\[
\begin{aligned}
\mathcal{P}_{\text{noise}}(y=0\mid s_1)
&= (1-\eta_{\text{err}})\cdot1 + \eta_{\text{err}}\cdot0 = 1-\eta_{\text{err}},\\
\mathcal{P}_{\text{noise}}(y=1\mid s_1)
&= (1-\eta_{\text{err}})\cdot0 + \eta_{\text{err}}\cdot1 = \eta_{\text{err}}.
\end{aligned}
\]
In World B (state $s_1$), $y=y^*$ directly:
\[
\mathcal{P}_{\text{hard}}(y=0\mid s_1)=1-\eta_{\text{amb}},\qquad
\mathcal{P}_{\text{hard}}(y=1\mid s_1)=\eta_{\text{amb}}.
\]
For state $s_2$, $\mathcal{P}(y=0\mid s_2)=1$,
$\mathcal{P}(y=1\mid s_2)=0$ in both worlds.  Hence
$\mathcal{P}(y\mid x)$ is identical across worlds.

\paragraph{Conditional expert distribution $\mathcal{P}(f_m\mid x)$.}
In World A (state $s_1$), experts predict the true label $y^*\equiv0$ with
accuracy $1-\varepsilon_1$:
\[
\mathcal{P}_{\text{noise}}(f_m=0\mid s_1)=1-\varepsilon_1,\qquad
\mathcal{P}_{\text{noise}}(f_m=1\mid s_1)=\varepsilon_1.
\]
In World B (state $s_1$):
\[
\begin{aligned}
\mathcal{P}_{\text{hard}}(f_m=0\mid s_1)
&= \mathbb{P}(y^*=0\mid s_1)\cdot\mathbb{P}(f_m=0\mid s_1,y^*=0)\\
&\quad + \mathbb{P}(y^*=1\mid s_1)\cdot\mathbb{P}(f_m=0\mid s_1,y^*=1)\\
&= (1-\eta_{\text{amb}})(1-\varepsilon_1) + \eta_{\text{amb}}(1-\varepsilon_1) = 1-\varepsilon_1,\\[4pt]
\mathcal{P}_{\text{hard}}(f_m=1\mid s_1)
&= (1-\eta_{\text{amb}})\varepsilon_1 + \eta_{\text{amb}}\varepsilon_1 = \varepsilon_1.
\end{aligned}
\]
For state $s_2$, both worlds give $\mathcal{P}(f_m=0\mid s_2)=1-\varepsilon_2$,
$\mathcal{P}(f_m=1\mid s_2)=\varepsilon_2$.

\paragraph{Joint identity.}
Combining the three components,
\[
\mathcal{P}_{\text{noise}}(x,y,\{f_m\})
= \mathcal{P}_{\text{hard}}(x,y,\{f_m\}),
\qquad \forall\,(x,y,\{f_m\}),
\]
which proves part~(iii) of Theorem~3. \hfill$\square$

\subsubsection{Algorithm Impossibility}
\label{sec:thm3_impossibility}

Let $\mathcal{A}$ be any algorithm that uses $n$ i.i.d.\ observations to
output per-sample noise flags $\{z_i\}_{i=1}^n$, where $z_i=1$ means ``label
is noisy.''  Since the observable distribution is identical in both worlds,
the algorithm's output distribution is also identical:
\[
\mathcal{L}_{\mathcal{P}_{\text{noise}}}(\mathcal{A}(D_n))
= \mathcal{L}_{\mathcal{P}_{\text{hard}}}(\mathcal{A}(D_n)).
\]

However, the \emph{correct} noise flags differ between worlds.  In World~A,
the truly noisy samples are those in state $s_1$ with $y=1$ (a proportion
$\rho\eta$ of the data).  In World~B, \emph{no} samples are noisy; the same
subset consists of correct but difficult samples.  The algorithm therefore
cannot simultaneously achieve zero error in both worlds.

Let $a$ be the expected fraction of the ambiguous subset
$\{x\in s_1,\;y=1\}$ that $\mathcal{A}$ flags as noisy.  Then:
\[
\mathrm{Error}_{\mathcal{P}_{\text{noise}}}(\mathcal{A}) = \rho\eta(1-a),\qquad
\mathrm{Error}_{\mathcal{P}_{\text{hard}}}(\mathcal{A}) = \rho\eta a.
\]
Hence
\[
\max\bigl(\mathrm{Error}_{\mathcal{P}_{\text{noise}}},\,
         \mathrm{Error}_{\mathcal{P}_{\text{hard}}}\bigr)
\ge \frac{\mathrm{Error}_{\text{noise}}+\mathrm{Error}_{\text{hard}}}{2}
= \frac{\rho\eta}{2}.
\]
This proves part~(iv). \hfill$\square$

\subsection{Extension to $K>2$ (Random Expert Construction)}
\label{sec:thm3_Kgeneral}

For $K>2$, the binary construction does not generalise directly because the
expert marginals in the two worlds no longer match automatically.  We instead
use a different construction where experts in the hardness world are
\emph{fully random} (conditionally independent of the true label).

\paragraph{World A (Noise).}
State $s_1$: $y^*\equiv0$.  Noise flips the label uniformly over the
remaining $K-1$ classes with probability $\eta$:
\[
\mathbb{P}(y=c\mid\text{noise},x\in s_1)=\frac{1}{K-1},\quad c\neq0.
\]
Expert $f_m$ predicts $0$ with probability $1-\varepsilon_1$, and each other
class with probability $\varepsilon_1/(K-1)$.  Thus
\[
\mathcal{P}_{\text{noise}}(y=0\mid s_1)=1-\eta,\quad
\mathcal{P}_{\text{noise}}(y=c\mid s_1)=\frac{\eta}{K-1}\;(c\neq0),
\]
\[
\mathcal{P}_{\text{noise}}(f_m=0\mid s_1)=1-\varepsilon_1,\quad
\mathcal{P}_{\text{noise}}(f_m=c\mid s_1)=\frac{\varepsilon_1}{K-1}\;(c\neq0).
\]
By Assumption~A4, noise and expert predictions are independent given the
state, so $y\perp f_m\mid s_1$ in World~A.

\paragraph{World B (Hardness).}
State $s_1$: all labels are correct ($y=y^*$).  The true label distribution
matches the World~A observed distribution:
\[
\mathbb{P}(y^*=0\mid s_1)=1-\eta,\quad
\mathbb{P}(y^*=c\mid s_1)=\frac{\eta}{K-1}\;(c\neq0).
\]
\textbf{Critical construction:} experts are fully random, independent of the
true label:
\[
\mathbb{P}(f_m=0\mid s_1)=1-\varepsilon_1,\quad
\mathbb{P}(f_m=c\mid s_1)=\frac{\varepsilon_1}{K-1}\;(c\neq0),
\]
with $f_m\perp y^*\mid s_1$ and $f_m\perp y\mid s_1$ (since $y=y^*$).

\paragraph{Equivalence.}
The marginal distributions of $y$, $f_m$, and the joint
$(y,\{f_m\})$ are identical in both worlds because both $y$ and $\{f_m\}$
are independent products of the same marginal distributions.  Hence for any
$K\ge2$:
\[
\mathcal{P}_{\text{noise}}(x,y,\{f_m\})
= \mathcal{P}_{\text{hard}}(x,y,\{f_m\}).
\]
The construction is extreme---experts in World~B are equivalent to random
guessers---but it suffices as an existence proof: there exists at least one
``hardness'' interpretation observationally indistinguishable from the
``noise'' interpretation. \hfill$\square$

\subsection{Corollaries: Which Assumptions Break Unidentifiability}
\label{sec:thm3_corollaries}

Theorem~3 shows that without additional assumptions, noise and difficulty are
indistinguishable.  Each corollary below identifies a \emph{sufficient
condition} that, if known to hold, breaks the equivalence.

\subsubsection{Anchor Labels (Corollary 1)}
\label{sec:corollary1}

\begin{corollary}[Anchor Labels Break Unidentifiability]
  If for a known subset $\mathcal{X}_{\text{anchor}}\subset\mathcal{X}$ the
  researcher has access to the true labels $y^*$ (e.g.\ via human review),
  then noise and difficulty are distinguishable on
  $\mathcal{X}_{\text{anchor}}$.  If $\mathcal{X}_{\text{anchor}}$ covers all
  states, the problem is globally identifiable.
\end{corollary}

\begin{proof}[Sketch]
  On anchor points, compare the observed label $y$ with the true label $y^*$.
  If $y\neq y^*$, the sample is necessarily noisy.  If $y=y^*$ but
  $f_m(x)\neq y$, the sample is necessarily difficult (the expert errs on a
  correctly labelled sample).  Extrapolating expert error rates from anchor
  points to non-anchor points breaks global unidentifiability.
\end{proof}

This corresponds to the SCX cold-start protocol, where a small number of
human-reviewed samples serve as anchors for initialising state-level noise
rate estimates.

\subsubsection{Disjoint Training Data (Corollary 2)}
\label{sec:corollary2}

\begin{corollary}[Known Disjoint Training Data Break Unidentifiability]
  If $M$ experts are trained on $M$ disjoint data subsets (Assumption~A1)
  and the researcher knows the training set composition, then noise and
  difficulty are distinguishable.
\end{corollary}

\begin{proof}[Sketch]
  Disjoint training implies expert errors are conditionally independent given
  $x$.  On noisy samples, \emph{all} experts have elevated error rates
  simultaneously (because they all predict the true label $y^*$, while the
  observed label $y$ differs).  On difficult samples, expert errors are
  independent.  The average error rate
  $\bar{e}(x)=\frac1M\sum_m\mathbf{1}\{f_m(x)\neq y\}$ satisfies
  \[
  \mathbb{E}[\bar{e}(x)\mid\text{noise}]
  = 1-\frac{1-\mathbb{E}[\bar{e}(x)\mid\text{clean}]}{K-1}
  \]
  (Theorem~1, Lemma~1).  This relation cannot hold simultaneously in a pure
  hardness world, providing a statistical test.
\end{proof}

\subsubsection{Input-Independent Noise (Corollary 3)}
\label{sec:corollary3}

\begin{corollary}[Input-Independent Noise Rate Breaks Unidentifiability]
  If the label noise rate $\eta(x)=\mathbb{P}(Y\neq Y^*\mid X=x)$ is
  constant across $x$ (Assumption~A4), and the expert error rates differ
  between states, then noise and difficulty are distinguishable.
\end{corollary}

\begin{proof}[Sketch]
  When $\eta(x)\equiv\eta$ is constant, the correlation pattern between the
  observed label~$y$ and expert errors~$\{f_m\}$ differs fundamentally
  between the two worlds.  In the noise world, $\mathbb{E}[\bar{e}(x)]$ is
  constant across states (marginalising over the label-flip event).  In the
  hardness world, $\bar{e}(x)$ varies with the state-dependent true label
  distribution.  Testing for input dependence of $\bar{e}(x)$ conditioned on
  $y$ distinguishes the two.
\end{proof}

\subsubsection{State Homogeneity (Corollary 4)}
\label{sec:corollary4}

\begin{corollary}[State Homogeneity Breaks Unidentifiability]
  If within each state $s$ the expected expert error rate on clean samples is
  constant (Assumption~A5), and states differ in this rate, then noise and
  difficulty are distinguishable.
\end{corollary}

\begin{proof}[Sketch]
  In the hardness world construction (World~B), state $s_1$ contains two
  subpopulations---samples with $y^*=0$ (expert accuracy $1-\varepsilon_1$)
  and samples with $y^*=1$ (expert accuracy $\varepsilon_1$).  This violates
  state homogeneity, which requires a single per-state error rate.
  Therefore, if Assumption~A5 holds, World~B's construction is invalid, and
  the unidentifiability is broken.
\end{proof}

\subsubsection{Balanced Error Distribution (Corollary 5)}
\label{sec:corollary5}

\begin{corollary}[Balanced Error Distribution Breaks Unidentifiability]
  If expert errors are uniformly distributed across all error classes
  (Assumption~A6 with $C_{\text{bal}}=1$), then noise and difficulty are
  distinguishable.
\end{corollary}

\begin{proof}[Sketch]
  In World~B, experts are biased toward class~0: when $y^*=1$, they predict
  class~0 with probability $1-\varepsilon_1$.  This creates a highly
  imbalanced error pattern (errors are almost exclusively in the
  $1\to0$ direction), violating the balanced error assumption.  If
  Assumption~A6 is known to hold, World~B is excluded.
\end{proof}

\subsection{Minimal Sufficient Assumption Set}
\label{sec:thm3_minimal}

Theorem~3's construction relies on specific conditions.  To exclude the
construction---and thereby break unidentifiability---at least one of the
following combinations of assumptions is necessary and sufficient:

\[
\begin{aligned}
\mathcal{A}_{\min}^{(1)} &= \{A1,\;A4,\;A5\},\\
\mathcal{A}_{\min}^{(2)} &= \{A1,\;A4,\;A6\},\\
\mathcal{A}_{\min}^{(3)} &= \{A5,\;A6\}\quad\text{with }|\mathcal{S}|\ge2.
\end{aligned}
\]

The SCX framework operates under $\mathcal{A}_{\min}^{(1)}$ (the standard
set A1--A6 provides additional technical benefits for finite-sample bounds).
Each assumption plays a distinct role:
\begin{itemize}
  \item \textbf{A1} (disjoint training): ensures expert independence, enabling
    Corollary~2's statistical test.
  \item \textbf{A4} (uniform noise): removes the possibility of
    input-dependent noise mimicking difficulty (Corollary~3).
  \item \textbf{A5} (state homogeneity): prevents a single state from
    containing both easy and difficult subpopulations (Corollary~4).
  \item \textbf{A6} (balanced errors): prevents expert bias from creating
    pseudo-difficulty patterns (Corollary~5).
\end{itemize}

If \emph{all} of these assumptions are violated, Theorem~3 applies fully:
no algorithm can distinguish label noise from sample difficulty.

% ===================================================================
\subsection{Corollary: The Everyone Equal Theorem (Everyone Equal Theorem)}
\label{sec:thm3_equal}

\begin{corollary}[The Everyone Equal Theorem (Everyone Equal Theorem)]
\label{cor:everyone-equal}
Let any observer $O$ claim to distinguish label noise from sample difficulty
using any algorithm $\mathcal{A}$. Then either
\begin{enumerate}[label=(\roman*)]
  \item the structural assumptions relied upon by $\mathcal{A}$ are
    explicitly stated and independently verifiable by any other observer,
    or
  \item the claim is mathematically empty---no amount of observational
    data can validate it, and any observer may legitimately reject it
    without further argument.
\end{enumerate}
Furthermore, if (i) holds, the assumptions are verifiable by any observer
without privileged access to $O$'s data, model, or expertise. All observers
are placed on equal epistemic footing.
\end{corollary}

\begin{proof}
Theorem~3 establishes $P_{W_A}(X,Y) = P_{W_B}(X,Y)$ for all observers,
i.e., the joint distribution of inputs, labels, and multi-expert predictions
is identical in the noise world and the difficulty world. Consequently, for
any discrimination rule $d$,
\[
\max_{W\in\{W_A,W_B\}} P_W\bigl(d(X,Y) \neq W\bigr) \geq \frac{1}{2},
\]
as proved via Le Cam's lemma in the main text. This lower bound is
observer-independent: it holds for \emph{any} algorithm $\mathcal{A}$
deployed by \emph{any} observer $O$, because the observational data do not
encode observer identity.

Therefore, if $O$ claims to achieve discrimination performance better than
chance, the claim must rely on assumptions that break the observational
equivalence $P_{W_A} = P_{W_B}$. These assumptions---the minimal sufficient
set $\mathcal{A}_{\min}$ derived above---must be stated. Once stated, they
are verifiable by any observer: checking disjoint training (A1), uniform
noise (A4), state homogeneity (A5), and balanced errors (A6) requires only
access to the training procedure and data distribution, not to $O$'s
internal state.

Hence no observer occupies a privileged epistemic position. The theorem
places all observers---regardless of institutional authority, computational
resources, or social status---on identical mathematical ground.
\end{proof}

\begin{remark}[No privileged observer]
The name ``The Everyone Equal Theorem'' (Everyone Equal Theorem) reflects the
corollary's operational meaning: the mathematical structure of data quality
assessment does not permit any person, algorithm, institution, or nation to
claim inherent authority over the distinction between noise and difficulty.
Authority must be earned through transparently stated, independently
verifiable assumptions. This is the formal basis for the SCX framework's
commitment to epistemic equality.
\end{remark}

% ===================================================================
\subsection{Conjecture: The Good Person Convergence Conjecture (Good Person Convergence)}
\label{sec:good-convergence}

\begin{conjecture}[The Good Person Convergence Conjecture (Good Person Convergence Conjecture)]
\label{conj:good-convergence}
Consider a community of $N$ agents interacting under the SCX framework:
each agent may produce data, claim data quality, or audit others' claims.
Let the following conditions hold:
\begin{enumerate}[label=(C\arabic*)]
  \item \textbf{The Honest Person Condition}: every data-quality claim must state its
    structural assumptions (Theorem~3);
  \item \textbf{The Everyone Equal Condition}: no agent occupies a privileged epistemic
    position (Corollary~\ref{cor:everyone-equal});
  \item \textbf{The No-Pretension Condition}: audit outputs are unconditionally verifiable
    by any third party (non-proliferation equilibrium, Theorem~NPE).
\end{enumerate}
Then, as $t\to\infty$:
\begin{enumerate}[label=(\roman*)]
  \item The fraction of agents making unverifiable claims converges to
    zero---``evil-doers bow'': dishonest behaviour becomes detectable and
    economically irrational;
  \item The expected payoff of honest behaviour exceeds that of dishonest
    behaviour for all agents---``the honest stand tall'': honesty is the strictly
    dominant strategy;
  \item The community converges to an equilibrium in which every agent
    simultaneously satisfies the three conditions, and no agent can improve
    their position by deviating.
\end{enumerate}
\end{conjecture}

\begin{remark}[Status of this claim]
\label{remark:good-convergence-status}
\textbf{This is a conjecture, not a theorem.} The original manuscript
presented a 13-line ``proof sketch'' that merely cited Theorems~3, NPE,
and Corollary~\ref{cor:everyone-equal} without establishing a formal
logical chain from premises to conclusion. A rigorous proof would require:
\begin{enumerate}[label=(\alph*)]
  \item A formal agent model (type space, strategy space, payoff functions);
  \item A dynamic process description (discrete/continuous time, update rules);
  \item A rigorous definition of convergence (topological space, convergence mode);
  \item A Lyapunov or potential function establishing the stability of the
    all-honest equilibrium;
  \item Verification that the three conditions (C1--C3) are not circular---each
    is itself a consequence of other SCX theorems, so their joint satisfaction
    must be proved consistent.
\end{enumerate}
Until such a proof is provided, the ``Good Person Convergence'' claim remains
a compelling research programme rather than an established result. The
intuition---that SCX's mathematical structure makes honesty strategically
dominant---is provocative and merits further investigation, but it has not
been formally demonstrated.
\end{remark}

% ===================================================================
\subsection{Epistemic Formalization System (E1-E5 + K Operator)}
\label{sec:epistemic}

We formalize the classical epistemological question---``when does a
community know a claim?''---as a probabilistic decision problem within
the SCX framework. The key insight: \emph{knowledge is a claim that
survives multi-expert audit with provable statistical guarantees.}

% -------------------------------------------------------------------
\subsubsection{Axioms for the Verification Framework}
\label{sec:epistemic-axioms}

\begin{definition}[Verification framework]
\label{def:verification-framework}
Let $\mathcal{S}$ be a claim with truth value $T(\mathcal{S})\in\{0,1\}$
(unknown to verifiers). Let $\mathcal{C}=\{v_1,\ldots,v_M\}$ be $M$
verifiers. Each $v_m$ produces a noisy binary judgment
$J_m(\mathcal{S})\in\{0,1\}$ via access only to observational data
$(X,Y,\{f_m\}_{m=1}^M)$ as defined in the SCX framework.
\end{definition}

\begin{assumption}[Epistemic axioms E1--E5]
\label{ax:epistemic}
\begin{enumerate}[label=\textbf{E\arabic*}, leftmargin=*, itemsep=4pt]
  \item \textbf{Assumption transparency:} the claim $\mathcal{S}$ is
    accompanied by an explicit assumption set $\mathcal{A}_{\mathcal{S}}$
    (enforced by Theorem~3, The Honest Person Theorem). If $\mathcal{A}_{\mathcal{S}}$
    is empty or unstated, $\mathcal{S}$ is rejected without evaluation.
  \item \textbf{Assumption soundness:} under $\mathcal{A}_{\mathcal{S}}$,
    verifier judgments satisfy
    $\mathbb{P}(J_m = T(\mathcal{S}) \mid \mathcal{A}_{\mathcal{S}}) = p_m > 1/2$.
    That is, each verifier is better than chance when assumptions hold.
    The verifier errors are conditionally independent given
    $\mathcal{A}_{\mathcal{S}}$ and $T(\mathcal{S})$.
  \item \textbf{State-homogeneous verifier quality:} within each state
    $s\in\mathcal{S}$, $\mathbb{P}(J_m \neq T(\mathcal{S}) \mid s) \leq
    \varepsilon_s < 1/2$, with $\Delta_s = 1/2 - \varepsilon_s > 0$.
  \item \textbf{Observational indistinguishability:} verifiers have access
    only to $(X,Y,\{f_m\})$---the same data as any SCX observer. They
    possess no oracle access to $T(\mathcal{S})$.
  \item \textbf{Counterexample logging:} the Yajie CEC $\mathcal{E}_t$
    accumulates all adjudicated judgments, so that any counterexample
    detected at time $t$ is permanently recorded.
\end{enumerate}
\end{assumption}

% -------------------------------------------------------------------
\subsubsection{The Knowledge Operator}
\label{sec:knowledge-operator}

\begin{definition}[SCX-knowledge operator $\mathbf{K}$]
\label{def:K-operator}
For claim $\mathcal{S}$, assumption set $\mathcal{A}_{\mathcal{S}}$,
community $\mathcal{C}$ of size $M$, and state $s\in\mathcal{S}$, define
the \textbf{verification score}:
\[
V_M(\mathcal{S}, s) = \frac{1}{M}\sum_{m=1}^{M} J_m(\mathcal{S}; s)
\]
where $J_m(\mathcal{S}; s)$ is verifier $m$'s judgment on samples in
state $s$. Define the \textbf{decision threshold} $\theta_s \in (0,1)$.

A \textbf{counterexample} is a triple $(x, y, \{f_m(x)\})$ such that
$\mathcal{A}_{\mathcal{S}}$ holds but $V_M(\mathcal{S}, s(x)) < \theta_{s(x)}$.
Let $\mathcal{E}_t$ be the set of counterexamples accumulated up to time $t$.
A counterexample $e \in \mathcal{E}_t$ is \textbf{resolved} if either (a) it
has been adjudicated and found consistent with $\mathcal{A}_{\mathcal{S}}$
under additional scrutiny, or (b) the claim is amended to accommodate it.
Otherwise it is \textbf{unresolved}.

The \textbf{SCX-knowledge operator} is defined per-state:
\[
\mathbf{K}_{t}(\mathcal{S}, \mathcal{C}, s) =
\begin{cases}
1 & \text{if } V_M(\mathcal{S}, s) \geq \theta_s
      \text{ and no unresolved counterexample in } \mathcal{E}_t
      \text{ involves state } s,\\[4pt]
0 & \text{otherwise.}
\end{cases}
\]
The claim-level knowledge operator aggregates state-level judgments:
$\mathbf{K}_{t}(\mathcal{S}, \mathcal{C}) = \mathbf{1}\{\mathbf{K}_t(\mathcal{S},
\mathcal{C}, s) = 1 \text{ for all } s \in \supp(
ho)\}$.
\end{definition}

% -------------------------------------------------------------------
\subsubsection{Rigorous Guarantee}
\label{sec:epistemic-guarantee}

\begin{theorem}[Epistemic Formalization System — E1-E5 + K Operator]
\label{thm:epistemic-rigorous}
Under the \textbf{SCX Axiom System (SCX Axiom System)} E1--E5 and the SCX framework conditions A1--A6
(Theorem~1), for any claim $\mathcal{S}$ with stated assumptions
$\mathcal{A}_{\mathcal{S}}$ and any community $\mathcal{C}$ of $M$
independent verifiers, the knowledge operator satisfies:
\[
\boxed{\mathbb{P}\bigl(\mathbf{K}(\mathcal{S},\mathcal{C}) \neq T(\mathcal{S})
  \mid \mathcal{A}_{\mathcal{S}}\bigr)
  \leq \sum_{s\in\mathcal{S}} \rho_s \exp\bigl(-2M\Delta_s^2\bigr)}
\]
where $\rho_s$ is the fraction of verification samples in state $s$,
and $\Delta_s = \theta_s - (1-p_s)$ is the separation between the
decision threshold and the expected error rate. The error probability
decays exponentially in the number of verifiers $M$.

Furthermore, this bound is tight: no knowledge operator operating under
E1--E5 can achieve a smaller error exponent without additional structural
assumptions, by Theorem~4 (The Extreme Precision Theorem).
\end{theorem}

\begin{proof}
The proof is a direct application of Theorem~1's concentration argument
to the verification setting. For each state $s$, the verifier judgments
$\{J_m(\mathcal{S}; s)\}_{m=1}^M$ are $M$ independent Bernoulli trials
with success probability $p_s$ (E2, E3). If $T(\mathcal{S})=1$ (the claim
is true), then $p_s \geq 1-\varepsilon_s$ by E3. The verification score
$V_M$ is a sample mean of $M$ independent Bernoulli variables.
By Hoeffding's inequality and the state-conditioned decomposition of
Theorem~1 (equation S1.7):
\[
\mathbb{P}\bigl(V_M \leq \theta_s \mid T(\mathcal{S})=1\bigr)
  \leq \exp\bigl(-2M(\theta_s - p_s)^2\bigr)
  = \exp\bigl(-2M\Delta_s^2\bigr).
\]
The same bound holds symmetrically when $T(\mathcal{S})=0$, with
$\Delta_s = p_s - (1-\theta_s)$.

The verifier errors are conditionally independent given the state (E3
strengthened by A2' from Theorem~1). The optional stopping argument from
Theorem~1's proof extends to the sequential audit setting, ensuring that
the bound does not degrade with multiple looks at the data.

For the tightness claim, Theorem~4 (The Extreme Precision Theorem) establishes that the
exponential rate $2M\Delta_s^2$ is minimax-optimal: any decision rule
operating on $M$ verifier judgments incurs an error exponent at least
this large. Hence the SCX-knowledge operator achieves the fundamental
limit. $\square$
\end{proof}

% -------------------------------------------------------------------
\subsubsection{Properties of $\mathbf{K}$}
\label{sec:K-properties}

\begin{corollary}[Knowledge operator properties]
\label{cor:K-properties}
The SCX-knowledge operator $\mathbf{K}$ satisfies:
\begin{enumerate}[label=(\roman*)]
  \item \textbf{Non-triviality:} $\mathbf{K}(\mathcal{S}) \neq 1$ for any
    $\mathcal{S}$ with $\mathcal{A}_{\mathcal{S}}=\emptyset$, because
    E1 rejects unstated assumptions;
  \item \textbf{Exponential reliability:}
    $\mathbb{P}(\mathbf{K} \text{ correct}) \to 1$ exponentially in $M$;
  \item \textbf{Monotonicity under audit:} if $\mathbf{K}_t(\mathcal{S})=1$,
    then $\mathbf{K}_{t+1}(\mathcal{S})=1$ unless a counterexample emerges
    (E5), in which case $\mathbf{K}$ flips to $0$ and stays there;
  \item \textbf{Gettier immunity (operational):} Gettier (1963) cases arise when a claim
    is justified but true only accidentally---the justification chain
    contains an unstated contingency. Under E1, any such contingency must
    appear in $\mathcal{A}_{\mathcal{S}}$, where it becomes either (a)
    verified (the claim holds under it) or (b) falsified (the claim
    fails). This addresses the \emph{operational} Gettier problem---cases
    where the contingency is expressible as a verifiable condition. It does
    \emph{not} claim to resolve \emph{all metaphysical} Gettier cases
    (e.g., those involving non-verifiable possible-worlds semantics);
  \item \textbf{Epistemic equality:} $\mathbf{K}$ depends only
    on $\mathcal{A}_{\mathcal{S}}$ and $\{J_m\}$, never on the identity,
    credentials, or institutional affiliation of the verifiers
    (Corollary~\ref{cor:everyone-equal}).
\end{enumerate}
\end{corollary}

% -------------------------------------------------------------------
\subsubsection{Connection to Classical Epistemology}
\label{sec:epistemic-classical}

\begin{remark}[What we did to Plato]
Plato (c.\ 369 BCE) defined knowledge as justified true belief.
Gettier (1963) found counterexamples where justification and truth
coincide accidentally due to unstated contingencies. The subsequent
literature attempted to \textit{redefine} justification.

The SCX approach does not redefine---it \textit{operationalizes}.
Knowledge is not a metaphysical state of an individual mind; it is
the mathematical status of a claim that has survived multi-expert
audit under stated assumptions with provable statistical guarantees.
The bound $\exp(-2M\Delta_s^2)$ replaces philosophical debate with
a quantitative threshold: ``how many independent verifiers do we
need?'' becomes a solvable equation.
\end{remark}
% ===================================================================

\subsection{Connections to Prior Art}
\label{sec:thm3_prior}

\subsubsection{Measurement Error Models}

In classical measurement error models
\cite{fuller1987measurement,carroll2006measurement}, the observed variable
$Y$ is modelled as $Y=Y^*+U$ where $U$ is measurement error.  Without
validation data or instrumental variables, the distributions of $Y^*$ and
$U$ are not separately identifiable from the marginal distribution of $Y$.
Theorem~3 extends this classical result from continuous measurement error to
the classification setting with multi-expert predictions.  The key novelty is
that even with the additional information $\{f_m(X)\}$---which one might
expect to resolve the ambiguity---the unidentifiability persists.

\subsubsection{Label Noise Transition Matrix Unidentifiability}

The label noise literature \cite{menon2015learning,patrini2017making}
established that the noise transition matrix $T_{ij}=\mathbb{P}(\tilde Y=j\mid
Y^*=i)$ is not identifiable from $(\tilde Y,X)$ alone without anchor points or
diagonal-dominance assumptions.  Theorem~3 generalises this to the multi-expert
setting, showing that even with $M$ expert predictions, the ambiguity between
noise and intrinsic difficulty remains.

\subsubsection{Dawid-Skene Unidentifiability}

The Dawid-Skene model \cite{dawid1979maximum} treats the true label as a latent
variable and each annotator's confusion matrix as a parameter.  Standard
identifiability issues include label-switching symmetries.  Theorem~3 addresses
a different and more fundamental ambiguity: even after resolving
label-switching, one cannot distinguish whether an annotator's error is due to
noise in the label or difficulty of the sample.  The two unidentifiability
problems are orthogonal: both must be resolved for reliable label quality
assessment.

\subsubsection{Causal Structure Unidentifiability}

From a causal perspective, Theorem~3 can be understood as a Markov equivalence
class problem \cite{pearl2009causality}.  Two causal directed acyclic graphs
generate identical observational distributions:
\begin{itemize}
  \item \textbf{Noise graph:} $S\to Y^*\to Y\leftarrow\text{Noise}$,
    $X\to\{f_m\}$, where Noise is an independent exogenous variable.
  \item \textbf{Difficulty graph:} $S\to Y^*(=Y)$, $X\to\{f_m\}$, where state
    $S$ affects both the true label distribution and the expert error rates.
\end{itemize}
These two graphs are observationally indistinguishable, forming a
Markov equivalence class of size~2.

\subsection*{References for Section~S3}

\begingroup
\small
\begin{thebibliography}{20}
\bibitem{fuller1987measurement} Fuller, W. A. \textit{Measurement Error Models}. Wiley (1987).
\bibitem{carroll2006measurement} Carroll, R. J., Ruppert, D., Stefanski, L. A. \& Crainiceanu, C. M. \textit{Measurement Error in Nonlinear Models: A Modern Perspective}. Chapman and Hall/CRC (2006).
\bibitem{menon2015learning} Menon, A. K., van Rooyen, B., Ong, C. S. \& Williamson, R. C. Learning from corrupted binary labels via class-probability estimation. \textit{Proc. ICML} (2015).
\bibitem{patrini2017making} Patrini, G., Rozza, A., Menon, A. K., Nock, R. \& Qu, L. Making deep neural networks robust to label noise: A loss correction approach. \textit{Proc. CVPR} (2017).
\bibitem{dawid1979maximum} Dawid, A. P. \& Skene, A. M. Maximum likelihood estimation of observer error-rates using the EM algorithm. \textit{Applied Statistics} \textbf{28}(1), 20--28 (1979).
\bibitem{pearl2009causality} Pearl, J. \textit{Causality: Models, Reasoning, and Inference}. Cambridge University Press, 2nd ed. (2009).
\bibitem{bahadur1960deviations} Bahadur, R. R. \& Rao, R. R. On deviations of the sample mean. \textit{Annals of Mathematical Statistics} \textbf{31}(4), 1015--1027 (1960).
\end{thebibliography}
\endgroup

% ============================================================
% Cross-Reference: Theorem~3 (The Honest Person Theorem (The Honest Person Theorem) / SCX's
% Uncertainty Principle) is treated in three companion documents:
%   - S3_thm3_unidentifiability.tex (this file): formal proof, minimum
%     sufficient assumption set, corollaries, connections to prior art.
%   - taxonomic_nn/theorem3.tex: standalone exposition emphasizing the
%     Uncertainty Principle framing and LLM hallucination application.
%   - SCX_HISTORY.tex (Sec "Theorem 3 as Uncertainty Principle"):
%     intellectual history, strategic framing, and implications for the
%     data-cleaning research programme.
% ============================================================

\endinput
